\scanpage{91}
\mspage{45}
$\mathcal E_U$ sur $\mathcal E_S$ ni $=$ de $\mathcal E_{\widetilde U}$ sur $\mathcal E_{\widetilde{\mathcal O}}$, définit néanmoins, \deleted{\unsure{leur|le|sur}} \unsure{le|de} sous-groupe $\mathcal E_{\widetilde K}=N_{s'}$ de $\mathcal E_K$ (qui est extension d'un $\mathcal E_k$ par un facteur trivial kummérien $T\underset{\deleted{\unsure{de|d}}}{\simeq}I_{s'}$, où $k$ est un corps de degré de \unsure{tr.-scend|tr.-scendance} $\deg\operatorname{tr}K-1$, lorsque $S$ est de type fini sur un corps de \unsure{réf|$N$}, disons $\mathbf Q$) \uline{presque} un relèvement --- seul le facteur $I_{s'}\simeq T$ y obstrue encore : en tous cas $N_{s'}\simeq\mathcal E_{\widetilde K}$ est plongé dans un $N_{\mathcal E_{X_s}}(L_i)$, et toute section de cette extension de $E_{k(s)}$ par $T$ donne donc une section de $\mathcal E_{X_s}$ sur $\mathcal E_{k(s)}$ --- mais de \uline{deuxième} espèce.

Une \deleted{\unsure{f}} autre façon de dire ces choses est la suivante. Une section rationnelle $f$ de $\frac{U}{S}$ sur $S$, de domaine de définition l'ouvert $S'$ de $S$, \deleted{\unsure{peut|permet}} peut s'interpréter en termes de la famille des pts $f(s)\in U_s$, où $s$ parcourt l'ens des pts de $S$
\margin{90}
